\begin{landscape}
\fontsize{9.5}{11.5}\selectfont
\setlength{\tabcolsep}{3pt}
\renewcommand{\arraystretch}{1.25}
\captionof{table}{Characteristics of Participants With and Without Subsequent Frailty Assessments.}\label{tab:landmark_fi_observability}
\vspace{0.5em}
\begin{tabular}{@{}L{4.4cm}*{3}{>{\raggedleft\arraybackslash}p{2.55cm}>{\raggedleft\arraybackslash}p{2.55cm}>{\raggedleft\arraybackslash}p{1.2cm}}@{}}
\toprule
Characteristic & \multicolumn{3}{c}{HRS} & \multicolumn{3}{c}{ELSA} & \multicolumn{3}{c}{KLoSA} \\
\cmidrule(lr){2-4}\cmidrule(lr){5-7}\cmidrule(lr){8-10}
 & \multicolumn{1}{c}{Observed FI} & \multicolumn{1}{c}{No observed FI} & \multicolumn{1}{c}{SMD} & \multicolumn{1}{c}{Observed FI} & \multicolumn{1}{c}{No observed FI} & \multicolumn{1}{c}{SMD} & \multicolumn{1}{c}{Observed FI} & \multicolumn{1}{c}{No observed FI} & \multicolumn{1}{c}{SMD} \\
\midrule
Participants, n & 5,114 & 1,054 & -- & 699 & 180 & -- & 2,502 & 248 & -- \\
Age at shock interview, years & 71.3 (8.0) & 77.5 (9.7) & -0.693 & 71.3 (7.2) & 75.7 (9.2) & -0.536 & 71.1 (6.9) & 77.0 (8.3) & -0.786 \\
FI at shock interview & 0.253 (0.169) & 0.369 (0.211) & -0.604 & 0.222 (0.159) & 0.283 (0.174) & -0.364 & 0.132 (0.104) & 0.241 (0.218) & -0.638 \\
FI at landmark interview & 0.276 (0.180) & 0.438 (0.234) & -0.777 & 0.241 (0.165) & 0.318 (0.197) & -0.425 & 0.145 (0.121) & 0.287 (0.240) & -0.751 \\
Substantial wealth recovery, n (\%) & 1,703 (33.3) & 293 (27.8) & +0.120 & 271 (38.8) & 60 (33.3) & +0.113 & 943 (37.7) & 80 (32.3) & +0.114 \\
\bottomrule
\end{tabular}\par
\vspace{0.6em}
\noindent\parbox{\linewidth}{\raggedright\fontsize{9}{11}\selectfont\textit{Notes:} Values are presented as mean (SD), participant counts, or n (\%). The comparison included participants with classifiable wealth-recovery status, valid landmark FI and interview timing, and complete information on shock FI and prespecified covariates, irrespective of subsequent FI availability. Within each cohort, observed FI denotes at least one valid assessment after the landmark interview; no observed FI denotes none. Age and shock FI were measured at the index-shock interview; landmark FI was measured when wealth recovery was assessed. For continuous variables, standardised mean differences (SMDs) were calculated within each cohort as the mean among participants with subsequent FI data minus that among those without, divided by the square root of the average of the two sample variances. SMDs are displayed to three decimal places; negative values indicate lower values among participants with subsequent FI data. Substantial wealth recovery was defined as regaining at least 50\% of the wealth lost during the index shock by the landmark interview. Percentages use the corresponding cohort × FI-availability group as the denominator. For this binary variable, the signed SMD was calculated as the recovery proportion in the observed-FI group minus that in the no-observed-FI group, divided by the square root of their average Bernoulli variance. Sample counts differ from the sequential selection process in Supplementary Figure S1 because eligibility was assessed independently of subsequent FI availability. Missing subsequent FI could reflect death, non-response, limited follow-up opportunity, unavailable FI data, or unknown reasons. These descriptive comparisons do not provide complete attribution of missingness or quantify selection bias. FI, frailty index; SMD, standardised mean difference; HRS, Health and Retirement Study; ELSA, English Longitudinal Study of Ageing; KLoSA, Korean Longitudinal Study of Aging.}
\end{landscape}
